% Drafted last; ~250 words; structured per NEJM AI envelope.
\begin{abstract}
\noindent\textbf{Background.} AI clinical decision support (CDSS)
tools are proliferating in practice, yet external validation remains
concentrated on a handful of public ICU electronic health record
datasets. Three problems result: (i)~\emph{benchmark
saturation}, repeated tuning on the same public test sets inflates
reported performance; (ii)~a \emph{reproducibility gap}, few
publications release runnable code, so retrospective re-execution at a
new site is slow or infeasible; and (iii)~\emph{incomplete
reporting}, most studies omit full TRIPOD+AI items and subgroup
fairness metrics that regulators increasingly require. Bilateral data
use agreements stretch multi-site validation over years and rarely
reach more than a few sites.

\noindent\textbf{Methods.} FLAIR is a federated benchmark in which
code travels to data, never the reverse. It is anchored on the Common
Longitudinal ICU Format (CLIF), a common data element--aligned
critical-care schema adopted by 12+ academic medical centers. Five
binary classification tasks span outcome, deterioration, and
intervention targets (ICU daily mortality, daily discharge to
long-term acute care, extubation failure, sepsis, unplanned ICU
readmission). Models run as committee-reviewed code or containers at each
site and emit aggregate-only JSON reports (cohort characteristics,
discrimination, calibration, and subgroup fairness metrics following
TRIPOD+AI) with HIPAA cell suppression. The full FLAIR stack, library, task
definitions, container/code specification, and report generators, is
released open source (Apache~2.0) at
\texttt{github.com/clif-consortium/FLAIR}, and the task catalog is
designed to grow: any contributor can add a new benchmark task as a
single-file Python module via GitHub pull request.

\noindent\textbf{Results.} Mode-matched baselines trained on MIMIC-CLIF
discriminate well on internal test (landmark mortality AUROC 0.86,
sepsis peak 0.69). At the first external site (Rush, 13{,}839 ICU
stays) discrimination drops on every task, from a modest decline on
mortality (0.86$\rightarrow$0.75) to a collapse on sepsis
(0.69$\rightarrow$0.41), quantifying the cross-site distribution shift
the benchmark exists to surface.

\noindent\textbf{Conclusions.} FLAIR shrinks retrospective multi-site
CDSS validation from years to weeks, exposes models to a growing
federated cohort that resists public-benchmark saturation, and
produces the TRIPOD+AI and fairness evidence package regulators
increasingly require, without patient-level data leaving any site.
The initial five-task suite is a seed, not a ceiling: the open
governance model invites the community to extend the benchmark to
new clinical questions over time.
\end{abstract}
